\documentclass[11pt]{article}
\usepackage[utf8]{inputenc}
\usepackage[T1]{fontenc}
\usepackage[spanish]{babel}
\usepackage{amsmath, amssymb, bm}
\usepackage{geometry}
\geometry{margin=2.5cm}

\title{Ridge/LS 05: relacion profunda con MMSE y estimacion de canal}
\author{}
\date{}

\begin{document}
\maketitle

\section{Punto de partida}

El operador:
\[
    (A^H A+\lambda I)^{-1}A^H
\]
aparece como solucion de un problema de minimos cuadrados regularizado. Pero
tambien puede entenderse como una forma simplificada de estimacion MMSE bajo
ciertas hipotesis estadisticas.

\section{Modelo estadistico}

Usamos:
\[
    \bm y = A\bm h + \bm n.
\]

Supongamos:
\[
    \bm h \sim \mathcal{CN}(\bm 0, R_{hh}),
    \qquad
    \bm n \sim \mathcal{CN}(\bm 0, R_{nn}).
\]

Aqui:

\begin{itemize}
    \item \(R_{hh}=E[\bm h\bm h^H]\) es la covarianza del canal.
    \item \(R_{nn}=E[\bm n\bm n^H]\) es la covarianza del ruido.
\end{itemize}

\section{Estimador LMMSE general}

El estimador lineal MMSE de \(\bm h\) a partir de \(\bm y\) es:
\[
    \hat{\bm h}_{LMMSE}
    =
    R_{hh}A^H
    \left(
        A R_{hh} A^H + R_{nn}
    \right)^{-1}
    \bm y.
\]

Esta expresion minimiza:
\[
    E\left[\|\bm h-\hat{\bm h}\|_2^2\right]
\]
dentro de la clase de estimadores lineales.

\section{Caso con covarianzas escalares}

Si asumimos:
\[
    R_{hh} = \sigma_h^2 I,
    \qquad
    R_{nn} = \sigma_n^2 I,
\]
entonces:
\[
    \hat{\bm h}_{LMMSE}
    =
    \sigma_h^2 A^H
    \left(
        \sigma_h^2 A A^H + \sigma_n^2 I
    \right)^{-1}
    \bm y.
\]

Existe una identidad util:
\[
    A^H(AA^H+\alpha I)^{-1}
    =
    (A^HA+\alpha I)^{-1}A^H.
\]

Con:
\[
    \alpha = \frac{\sigma_n^2}{\sigma_h^2},
\]
obtenemos:
\[
    \hat{\bm h}_{LMMSE}
    =
    (A^H A+\alpha I)^{-1}A^H\bm y.
\]

Por tanto:
\[
    \lambda = \frac{\sigma_n^2}{\sigma_h^2}.
\]

\section{Interpretacion de \(\lambda\) como relacion ruido/canal}

Si el ruido es grande:
\[
    \sigma_n^2 \uparrow
    \quad \Rightarrow \quad
    \lambda \uparrow.
\]

El estimador se vuelve mas conservador.

Si el canal esperado tiene mucha energia:
\[
    \sigma_h^2 \uparrow
    \quad \Rightarrow \quad
    \lambda \downarrow.
\]

El estimador confia mas en las observaciones.

En el codigo, sin embargo:
\[
    \lambda = 10^{-2}
\]
es fijo. No se calcula automaticamente desde \(\sigma_n^2\) y \(\sigma_h^2\).

\section{Conexion con MAP}

Tambien se puede obtener la misma forma con estimacion MAP.

Si:
\[
    \bm n \sim \mathcal{CN}(0,\sigma_n^2I),
\]
la verosimilitud es proporcional a:
\[
    \exp\left(
        -\frac{1}{\sigma_n^2}
        \|\bm y-A\bm h\|^2
    \right).
\]

Si:
\[
    \bm h \sim \mathcal{CN}(0,\sigma_h^2I),
\]
el prior es proporcional a:
\[
    \exp\left(
        -\frac{1}{\sigma_h^2}
        \|\bm h\|^2
    \right).
\]

Maximizar el posterior equivale a minimizar:
\[
    \frac{1}{\sigma_n^2}\|\bm y-A\bm h\|^2
    +
    \frac{1}{\sigma_h^2}\|\bm h\|^2.
\]

Multiplicando por \(\sigma_n^2\):
\[
    \|\bm y-A\bm h\|^2
    +
    \frac{\sigma_n^2}{\sigma_h^2}
    \|\bm h\|^2.
\]

De nuevo:
\[
    \lambda = \frac{\sigma_n^2}{\sigma_h^2}.
\]

\section{Version con covarianza no diagonal}

Un MMSE mas completo no asume:
\[
    R_{hh} = \sigma_h^2 I.
\]

En canales reales, los taps pueden tener potencias distintas:
\[
    R_{hh}
    =
    \mathrm{diag}
    (\sigma_{h,0}^2,\sigma_{h,1}^2,\ldots,\sigma_{h,L-1}^2).
\]

Esto significa que algunos retardos son mas probables o mas fuertes que otros.

Si \(R_{hh}\) es diagonal pero no escalar, el estimador penaliza taps de forma
diferente. La regularizacion ya no seria \(\lambda I\), sino una matriz.

\section{Regularizacion matricial}

El problema regularizado general puede escribirse como:
\[
    \min_{\bm h}
    \|\bm y-A\bm h\|^2
    +
    \bm h^H R \bm h.
\]

La solucion es:
\[
    \hat{\bm h}
    =
    (A^H A+R)^{-1}A^H\bm y.
\]

Si:
\[
    R=\lambda I,
\]
todos los taps se penalizan igual.

Si:
\[
    R=\mathrm{diag}(\lambda_0,\lambda_1,\ldots,\lambda_{L-1}),
\]
cada tap tiene una penalizacion propia.

\section{Eigenvalues y MMSE}

Usando SVD:
\[
    A=U\Sigma V^H.
\]

El estimador ridge filtra cada direccion con:
\[
    f_i(\lambda)
    =
    \frac{\sigma_i}{\sigma_i^2+\lambda}.
\]

Esto se parece al principio MMSE:

\begin{itemize}
    \item si una direccion esta bien observada, \(\sigma_i\) grande, se conserva;
    \item si una direccion esta mal observada, \(\sigma_i\) pequeno, se atenua;
    \item si el ruido esperado es grande, \(\lambda\) aumenta y se atenua mas.
\end{itemize}

\section{Diferencia con ecualizador MMSE}

Es importante separar dos usos de la palabra MMSE.

\subsection{MMSE para estimar canal}

Aqui se busca:
\[
    \hat{\bm h}
    \approx
    \bm h.
\]

La regularizacion del codigo pertenece a esta etapa.

\subsection{MMSE para ecualizar datos}

Despues de estimar el canal, se podria ecualizar con:
\[
    \hat{X}_{MMSE}[k]
    =
    \frac{\hat{H}^*[k]}
    {|\hat{H}[k]|^2+\sigma_w^2/\sigma_x^2}
    Y[k].
\]

El programa actualmente no hace esto. El programa hace:
\[
    \hat{X}_{ZF}[k]
    =
    \frac{Y[k]}{\hat{H}[k]}.
\]

Por tanto:

\[
\boxed{
\text{El codigo usa regularizacion tipo LMMSE en estimacion de canal,
pero ecualizacion final ZF.}
}
\]

\section{Aplicacion directa al simulador}

En el codigo:

\[
    A[m,\ell] = e^{-j2\pi k_m\ell/N}.
\]

El vector \(\bm y\) es:
\[
    \bm y =
    \left[
        \hat{H}_{LS}[k_1],
        \hat{H}_{LS}[k_2],
        \ldots,
        \hat{H}_{LS}[k_M]
    \right]^T.
\]

El estimador produce:
\[
    \hat{\bm h}
    =
    (A^HA+\lambda I)^{-1}A^H\bm y.
\]

Luego:
\[
    \hat{\bm H}
    =
    F\hat{\bm h}.
\]

Y finalmente:
\[
    \hat{\bm X}
    =
    \bm Y \oslash \hat{\bm H},
\]
donde \(\oslash\) representa division elemento a elemento.

\section{Por que esto ayuda con los pilotos}

Si solo se interpolan pilotos, el receptor no usa explicitamente que el canal
vive en pocos taps. Con ridge/DFT, si lo usa:

\[
    \bm h =
    [h[0],h[1],\ldots,h[L-1]]^T.
\]

Esto impone estructura fisica. El canal no es una curva arbitraria en
frecuencia; viene de una respuesta temporal corta.

\section{Resumen final}

La expresion:
\[
    (A^H A+\lambda I)^{-1}A^H
\]
puede verse como:

\begin{enumerate}
    \item una pseudoinversa regularizada;
    \item la solucion de minimos cuadrados con penalizacion;
    \item un filtro espectral que controla valores singulares pequenos;
    \item una forma LMMSE simplificada si \(\lambda=\sigma_n^2/\sigma_h^2\);
    \item una herramienta para estimar taps de canal desde pilotos.
\end{enumerate}

En el simulador, su papel exacto es:
\[
\boxed{
\text{convertir mediciones LS en pilotos en una estimacion estable de }
\bm h.
}
\]

\end{document}
